% Isolated proposed B173, checkpoint36. No new deletion search or producer.
\subsection*{A complementary actual target retaining the noncommuting contact}
\statement{B.173}{Proved}{The actual retained two-color quotient has an entire three-generator source image with full induced topology, an exact source normal kernel, and a continuous compact contact section complementary to the one-sided section}
\label{stmt:B173}
Use the actual B.132 index-two quotient. Thus
\[
 W^{(2)}=\ker\bigl(W\xrightarrow{\chi_\Sigma}\mathbb Z/2\bigr),\qquad
 P=W^{(2)}\rtimes\langle m\rangle
       \longrightarrow G_7=Q_7\rtimes_\alpha\langle m\rangle,
\]
where
\[
 Q_7=F(D_0,B)*T^{(2)},\quad
 T^{(2)}=\ker\bigl(F(f,g,x)\xrightarrow{\chi_\Sigma}\mathbb Z/2\bigr),
 \quad D_0=x\bar d x^{-1},\quad\kappa=f^{-1}g.
\]
The superscript denotes the actual index-two subgroup. All original
node sheets and signed translates are killed by this proved quotient;
it retains both the positive \(B\)-core and the barred \(d\)-core.
It is a map of \(P\), not an asserted map
\(\Gamma_0\to G_7\) or an \(S\)-equivariant free-fiber section.

The whole actual \(H_5\)-image is
\begin{equation}\label{eq:B173-image}
 H_2=\langle m,B,D_0\rangle
       \cong(\mathbb Z_m\times\mathbb Z_{D_0})*\mathbb Z_B.
\end{equation}
It has full induced profinite topology in \(G_7\).
The actual restriction \(H_5\to H_2\) extends to a continuous epimorphism
\[
 \pi_2:\mathcal H_5\twoheadrightarrow
 \mathcal H_2=
   (\widehat{\mathbb Z}_m\times\widehat{\mathbb Z}_{D_0})
       \amalg\widehat{\mathbb Z}_B,
\]
with kernel
\begin{equation}\label{eq:B173-source-kernel}
 \ker\pi_2=
 \overline{\langle\!\langle
     A_s,E_s,[m,B_sD_sB_s^{-1}]
                  \rangle\!\rangle}_{\mathcal H_5}.
\end{equation}
The ordinary restriction has the corresponding ordinary normal kernel.

Put
\[
 \mathcal V=\ker(\mathcal H_2\to\widehat{\mathbb Z}_m).
\]
Then \(\mathcal V\) is the free profinite group on the continuous compact
family consisting of one point \(D_0\) and
\(B_z=m^zBm^{-z}\), \(z\in\widehat{\mathbb Z}\).
One has \(\pi_2(\mathcal I_5)=\mathcal V\), with a continuous section
\begin{equation}\label{eq:B173-section}
 D_0\longmapsto B_sD_sB_s^{-1},\qquad
 B_z\longmapsto m^zB_sm^{-z}.
\end{equation}
It is an actual compact section, but is not \(m\)-equivariant.
Neither an original presentation of \(H_5\) nor the full
fixed-generation equality is inferred.

\begin{proof}
\emph{B.173.1. Exact retained source and operator data.}
B.132 proves the actual epimorphism
\(W^{(2)}\to Q_7\) by keeping the even-coset \(B,\bar d\) cores
and all five retained Schreier generators. The other fourteen family
generators are killed. Its all-word character proof kills every whole
node, including all signed \(v\)-translates, even though the resulting
free quotient is used with the \(m\)-action only.
The actual signed \(m\)-rows give
\[
 \alpha(D_0)=D_0,\qquad
 \alpha(B)=\kappa B\kappa^{-1},\qquad
 \alpha(T^{(2)})=T^{(2)}.
 \tag{B.173.1}
\]
On the larger retained free group their defining formulas are
\[
 \begin{gathered}
 \alpha(f)=f^{-1}gf,\quad \alpha(g)=f,\quad
 \alpha(x)=xg^{-1}f,\\
 \alpha^{-1}(f)=g,\quad\alpha^{-1}(g)=gfg^{-1},\quad
 \alpha^{-1}(x)=xfg^{-1}.
 \end{gathered}
\]
They preserve \(\chi_\Sigma\) and restrict to the literal inverse
automorphisms of \(T^{(2)}\). In the original raw contact order,
\[
 (q(s_1),q(s_2),q(s_3),q(s_4))
       =(B^{-1},1,D_0^{-1},1).
\]
Substitution in the corrected contact basis therefore gives
\begin{equation}\label{eq:B173-contact-images}
 A_s,E_s\longmapsto1,\qquad
 B_s\longmapsto B,\qquad D_s\longmapsto B^{-1}D_0B.
\end{equation}
In particular these images generate \(H_2\). This is the actual
retained-source computation, not a replacement of \(D_s\) by an
uncorrected norm.

\emph{B.173.2. Two cyclic edges and their full intersection.}
Let \(G_T=T^{(2)}\rtimes\langle m\rangle\) and
\(t=\kappa^{-1}m\). The full ordinary target has presentation
\begin{equation}\label{eq:B173-two-edges}
 G_7=G_T*_{\langle m\rangle}
        (\langle m\rangle\times\langle D_0\rangle)
       *_{\langle t\rangle}
        (\langle t\rangle\times\langle B\rangle).
\end{equation}
The new commuting relations are exactly
\([m,D_0]=1\) and \([t,B]=1\), which recover all the displayed
fiber-action formulas. The two presentation substitutions are inverse.

Characteristic finite fiber refinements and a retained \(m\)-order
factor give
\[
 \widehat{G_T}=\widehat{T^{(2)}}\rtimes\widehat{\mathbb Z}_m.
 \tag{B.173.2}
\]
Indeed they refine every finite target of \(G_T\), exactly as in
B.169.2. The group \(\widehat{T^{(2)}}\) is its actual open subgroup
inside \(\widehat F(f,g,x)\), because \(T^{(2)}\) has finite index.
Let \(M=\overline{\langle m\rangle}\), \(T=\overline{\langle t\rangle}\)
inside this completion. With \(q_0=gf\), \(q_1=fg=fq_0f^{-1}\),
the entire B.132 ordered formula is
\[
 \kappa_{2n}=q_0^{-n}q_1^n,\qquad
 \kappa_{2n+1}=q_0^{-n}(f^{-1}g)q_1^n,\qquad
 t^n=\kappa_n^{-1}m^n.
 \tag{B.173.3}
\]
It holds for every integer, including negative integers.
Define \(\kappa_z=m^zt^{-z}\); continuity gives the same formulas
for all profinite \(n\). An element of \(M\cap T\) forces
\(\kappa_z=1\) by its source projection. Odd \(z\) is excluded by
abelianization. For \(z=2w\), the equality is
\(q_0^{-w}fq_0^wf^{-1}=1\) inside \(\widehat F(f,g,x)\).
The full primitive centralizer of \(f\) is
\(\widehat{\mathbb Z}_f\); the \(g\)-exponent of \(q_0^w\) then forces
\(w=0\). Hence
\begin{equation}\label{eq:B173-cyclic-intersection}
                            M\cap T=1.
\end{equation}

The ordinary version proves \eqref{eq:B173-image} directly.
In the first amalgam in \eqref{eq:B173-two-edges}, the subgroup
\(\langle m,D_0\rangle=\mathbb Z^2\) meets \(G_T\) in
\(\langle m\rangle\), and hence meets \(\langle t\rangle\) trivially.
In the second amalgam, a reduced word alternating nontrivial elements
of this \(\mathbb Z^2\) with nonzero powers of \(B\) is therefore an
amalgam-reduced word. It is nontrivial. Thus their subgroup is the
free product in \eqref{eq:B173-image}.

\emph{B.173.3. Every finite image target and the whole-source refinement.}
Given any finite \(\psi:H_2\to K\), its \(m\)-restriction extends to
\(M\). If its kernel is \(M_0\), use
\eqref{eq:B173-cyclic-intersection} and compactness to choose an open
normal head kernel \(V\) with
\[
 V\cap M\le M_0,\qquad M\cap TV\le M_0.
\]
In the finite head \(D=\widehat{G_T}/V\), the map
\(\rho:\overline M\to K\) then kills
\(\overline M\cap\overline T\).
On \(\Omega=D\times_{\overline M}K\), with
\([gh,k]=[g,\rho(h)k]\), let \(D_0\) act on the identity fiber by
the regular permutation \(\psi(D_0)\) and trivially off that fiber.
It commutes with \(\overline M\). Extend the regular permutation
\(\psi(B)\) from that fiber over its \(\overline T\)-orbit by
\[
 B[t_0,k]=[t_0,\psi(B)k],
\]
and set it to the identity on the complement.
This is well-defined because the intersection acts trivially on the
fiber, is invertible by the analogous inverse formula, and commutes
with \(\overline T\). Thus every relation of
\eqref{eq:B173-two-edges} holds in a finite permutation target.
Its \(H_2\)-restriction on the invariant identity fiber is the regular
\(\psi\)-action. Its kernel therefore refines \(\ker\psi\), proving
full topology of \(H_2\) in \(G_7\).

One must also retain the entire original \(S\)-action.
Let \(\eta:G_7\to L\) be such a finite target, and pull it back to
the actual \(P=W^{(2)}\rtimes\langle m\rangle\).
Its kernel on \(W^{(2)}\) has finite index in \(W\).
Core it in \(W\), and then take a characteristic finite-index subgroup
\(C\) of the finitely generated \(W\) inside that core.
Every actual signed source automorphism preserves \(C\).
Let \(\ell\) be the order of the prescribed image of \(m\), and take
the finite image
\[
 A_S=\operatorname{im}\bigl(
 S\to\operatorname{Aut}(W/C)\times\mathbb Z/\ell\bigr),
 \quad m\mapsto(\overline\alpha_m,1),\quad
 v\mapsto(\overline\alpha_v,0).
\]
Then the actual whole-group map
\(\Gamma_0\to(W/C)\rtimes A_S\) refines the original target on \(P\):
if \(wm^n\) is killed, \(w\in C\) and \(\ell\mid n\), so its
\(\eta\)-image is one. This retains the full source kernel, not
just its fiber part. It does not assume that \(v\) acts on \(Q_7\).

For each finite quotient of \(H_2\), the corresponding restriction
of \(H_5\to H_2\) is therefore continuous for the actual ambient
topology on \(H_5\). It extends to \(\mathcal H_5\).
Uniqueness on dense \(H_5\) makes these finite extensions compatible,
giving \(\pi_2:\mathcal H_5\to\widehat H_2\).
Its compact image contains dense \(H_2\), so it is surjective.
This proves the asserted actual compact map and all its topology claims.

\emph{B.173.4. Exact source kernels.}
Use the marked epimorphism
\[
 P_0=\widehat{\mathbb Z}_{A_s,m}^{\,2}
                 \amalg\widehat F(B_s,D_s,E_s)
          \twoheadrightarrow\mathcal H_5.
\]
The composite to \(\widehat H_2\) kills
\(A_s,E_s,[m,B_sD_sB_s^{-1}]\).
Conversely these relations give its exact presentation: after deleting
\(A_s,E_s\), replace \(D_s\) by
\(D_0=B_sD_sB_s^{-1}\), with inverse
\(D_s=B_s^{-1}D_0B_s\). The remaining presentation is precisely
\([m,D_0]=1\), with \(B_s\) free. Both maps fix the three new generators.
The composite kernel in \(P_0\) is therefore the closed normal closure
of the three displayed words. Its image in \(\mathcal H_5\) proves
\eqref{eq:B173-source-kernel}; the unknown kernel of
\(P_0\to\mathcal H_5\) is merely contained in it.
The identical ordinary argument proves the ordinary normal-kernel claim.
Unlike \(\mathcal K_5\), this full kernel contains \(A_s\) and is not
being asserted to lie in \(\mathcal I_5\).

\emph{B.173.5. The continuous contact family and its literal section.}
For every positive \(n\), take the free profinite group on
\(D_0\) and \(B_i\), \(i\in\mathbb Z/n\), with reduction maps on
the labels. Their inverse limit is universal for continuous maps from
the disjoint union of one point and \(\widehat{\mathbb Z}\):
a finite target factors through a residue partition, and compatible
maps to all finite quotients give the general case.
The label family topologically generates this limit.
Let \(m\) fix \(D_0\) and translate the \(B\)-index.
This jointly continuous action defines a semidirect product with
\(\widehat{\mathbb Z}_m\). Its presentation maps to and from
\(\widehat H_2\) fix \(m,D_0,B_0\), with the reverse family map
\(B_z\mapsto m^zBm^{-z}\). They are inverse.
This proves the asserted description of \(\mathcal V\).

The source projection gives \(\pi_2(\mathcal I_5)\subseteq\mathcal V\).
Conversely \(B_sD_sB_s^{-1}\) and every
\(m^zB_sm^{-z}\) lie in \(\mathcal I_5\), and their images are the
whole compact generating family of \(\mathcal V\). Its compact image
is therefore all of \(\mathcal V\).
The continuous label maps \eqref{eq:B173-section} extend by this
universal property to a section; their composite with \(\pi_2\)
fixes the labels. The section is injective and its compact image has
its actual intrinsic topology.

\emph{B.173.6. The complementary defects and exact retained correction.}
Write \(r_D=[m,D_s]\) and \(r_B=[m,B_sD_sB_s^{-1}]\).
The target of B.169 kills \(r_D\), but its image of \(r_B\) is
nontrivial by the actual B.168 certificate.
The present target kills \(r_B\), while
\[
 \pi_2(r_D)=B_1^{-1}D_0B_1B_0^{-1}D_0^{-1}B_0\ne1.
 \tag{B.173.4}
\]
Nontriviality follows already by projecting the continuous family to
its finite residue grid modulo two: this is a reduced nonempty word
in the distinct free generators \(D_0,B_0,B_1\).
Thus the new section cannot be \(m\)-equivariant; its lift of \(D_0\)
does not commute with \(m\) in the actual source.

The new section image lies inside the earlier
\(\mathcal U_{\mathrm{lift}}\), since its generators are words in
\(D_s,B_s\) and the earlier \(B_{0,z}\)-family. Nevertheless the
new projection is nontrivial on the earlier kernel \(\mathcal K_5\).
Precisely, continuity and the normal-closure formula of B.170 give
\begin{equation}\label{eq:B173-retained-defect-image}
 \begin{split}
 \pi_2(\mathcal K_5)
   &=\overline{\langle\!\langle[m,B^{-1}D_0B]
                             \rangle\!\rangle}_{\mathcal H_2}\\
   &=\overline{\langle\!\langle
        B_{z+1}^{-1}D_0B_{z+1}B_z^{-1}D_0^{-1}B_z:
                     z\in\widehat{\mathbb Z}
                             \rangle\!\rangle}_{\mathcal V}.
 \end{split}
\end{equation}
The second equality retains every compact \(m\)-translate before
taking the normal closure in \(\mathcal V\).
This identifies a genuine retained image of the noncommuting
correction. It does not identify its kernel in the original source,
and does not prove the full contact-image equality.
\end{proof}
